\section{Introduction}
Self-improving agents convert interaction histories into persistent state. External memory is a common implementation because it can update quickly without modifying model parameters. ReasoningBank writes reasoning memories from completed trajectories \citep{ouyang2026reasoningbank}. Agent Workflow Memory stores reusable procedures from prior interactions \citep{wang2024awm}. These systems make experience reusable across future tasks.

The memory writer needs a learning signal. In the released ReasoningBank implementation used by the recent fragility study \citep{ye2026fragility}, the terminal score controls which reflection objective writes the memory. Score one activates a success-reflection prompt. Score zero activates a failure-reflection prompt. The success writer asks what made the trajectory work. The failure writer asks what caused the attempt to fail. A one-bit reward change therefore changes the interpretation assigned to the same action history.

This implementation creates a persistent reward channel. A proxy evaluator can make an error at the end of an episode. The error then selects the wrong memory-writing objective. The resulting text is stored and becomes available to later tasks. Reward noise has crossed the episode boundary and entered the agent state.

The same fragility study reports substantial variance across repeated self-improvement runs. Its main experiments use ground-truth reward specifically to remove reward noise from the analysis \citep{ye2026fragility}. This design choice exposes an open causal question: how much of self-improvement variance can originate from reward-conditioned memory writing itself? Memory Reward Inflation supplies a complementary observation. Proxy reward errors can accumulate under repeated memory reuse and distort future learning value \citep{asadolahi2026inflation}.

We test the first link directly using released trajectories. For each sampled trajectory, the action history remains identical. The task remains identical. The writer model remains identical. We change only the reward label that selects the ReasoningBank reflection mode. Four complete paired interventions all produce different memory text. Their title sets also differ in every pair. Mean token-set Jaccard distance is 0.735.

A fresh prompt-perturbation control tests whether this result reduces to generic prompt wording sensitivity. On eight trajectories excluded from F0, semantic-preserving same-mode paraphrases change the system prompt more strongly than the released reward-mode contrast under token-set Jaccard distance. Between reward modes, mean memory distance is 0.700. The average same-mode control distance is 0.595. Their paired difference is 0.105 with exact sign-flip $p=0.0078$. The reward-conditioned contrast therefore survives a stronger same-mode wording control.

The paired experiment and prompt control establish a write-channel mechanism. A reward-label perturbation can survive as a persistent memory perturbation even when the underlying trajectory is fixed. We then connect this mechanism to downstream variance with a matched memory-injection protocol. Paired memories are injected into the same future task under a fixed environment snapshot. Repeated runs measure early action divergence and terminal success dispersion.

Our contribution is a causal decomposition of reward noise in self-improving agents. Reward error first changes the memory-writing objective. The written memory then changes future context. Reuse can amplify this state perturbation across later episodes. This view turns evaluator reliability into a state-consistency requirement for systems that learn from persistent memory.
